\documentclass[a4paper,12pt]{article}
\usepackage[T2A]{fontenc}
\usepackage[utf8]{inputenc}
\usepackage[russian]{babel}
\usepackage{amsmath,amssymb,mathtools}
\usepackage{helvet}
\renewcommand{\sfdefault}{cmss}
\renewcommand{\familydefault}{\sfdefault}
\usepackage{icomma}
\usepackage{geometry}
\geometry{left=22mm,right=22mm,top=22mm,bottom=23mm}
\usepackage{microtype}
\usepackage{tikz}
\usetikzlibrary{arrows.meta,positioning,shapes.geometric,calc}
\usepackage{tabularx,booktabs,array}
\usepackage{enumitem}
\usepackage{hyperref}
\usepackage{parskip}
\usepackage{fancyhdr}
\usepackage{xcolor}
\usepackage{setspace}
\setstretch{1.24}
\definecolor{accent}{HTML}{008080}
\definecolor{darkaccent}{HTML}{00585B}
\definecolor{ink}{HTML}{293033}
\definecolor{hint}{HTML}{5A676B}
\hypersetup{colorlinks=true,linkcolor=darkaccent,urlcolor=darkaccent}
\setlength{\parindent}{0pt}
\setlength{\emergencystretch}{3em}
\setlist[itemize]{leftmargin=5.5mm,itemsep=1pt,topsep=3pt}
\setlist[enumerate]{leftmargin=7mm,itemsep=2pt,topsep=3pt}
\renewcommand{\arraystretch}{1.28}
\pagestyle{fancy}
\fancyhf{}
\fancyhead[L]{\footnotesize\color{hint} Движение протяжённых тел и совместная работа}
\fancyhead[R]{\footnotesize\color{hint} Анна Баннова}
\fancyfoot[C]{\footnotesize\color{hint}\thepage}
\renewcommand{\headrulewidth}{0.25pt}
\renewcommand{\footrulewidth}{0pt}
\newcommand{\key}[1]{\textcolor{darkaccent}{\textbf{#1}}}
\newcommand{\ans}[1]{\textcolor{darkaccent}{\textbf{#1}}}
\newcommand{\ex}[1]{\par\smallskip\noindent\key{Пример.} \textbf{#1}\par\smallskip}
\newcommand{\task}[2]{\par\smallskip\noindent\textbf{№ #1.}\enspace #2\par}
\newcolumntype{Y}{>{\raggedright\arraybackslash}X}
\begin{document}
\thispagestyle{empty}
\begin{center}
\vspace*{24mm}
{\Large\color{darkaccent} МАТЕМАТИКА}\par
\vspace{20mm}
{\fontsize{29}{34}\selectfont\color{ink} Чек-лист}\par
\vspace{11mm}
{\fontsize{20}{27}\selectfont\color{accent} Движение протяжённых тел\\[2mm] и задачи на совместную работу}\par
\vspace{18mm}
{\large Подготовка к ЕГЭ по математике}\par
\vfill
{\Large Анна Баннова}\par
\vspace{11mm}
{\large 15 ноября 2026 г.}\par
\vspace{17mm}
{\normalsize Лёвин Артём Александрович\\эксклюзивно для Анны Банновой}\par
\vspace{14mm}
\begin{tikzpicture}[x=1cm,y=1cm]
\draw[darkaccent,line width=0.9pt] (-5.2,0) .. controls (-3.4,0.9) and (-1.6,-0.7) .. (0,0)
 .. controls (1.5,0.7) and (3.4,-0.8) .. (5.2,0);
\fill[darkaccent] (0,0) circle (1.7pt);
\end{tikzpicture}
\end{center}
\clearpage

\section*{1. Когда длина движущегося тела имеет значение}
У поезда, баржи и других протяжённых объектов есть собственная длина. Чтобы объект \emph{полностью прошёл} мимо наблюдателя, моста или другого поезда, нужно учитывать расстояние, которое проходит его передняя часть относительно выбранного ориентира.

\key{Главная модель:} время полного прохождения равно относительному пути, делённому на относительную скорость:
\[
\boxed{\displaystyle t=\frac{S_{\text{отн}}}{v_{\text{отн}}}},\qquad
S_{\text{отн}}=v_{\text{отн}}t.
\]
В каждом задании нужно установить две вещи: \emph{что именно проходит мимо чего} и \emph{как объекты движутся относительно друг друга}.

\section*{2. Относительная скорость}
Обозначим скорости объектов через $v_1$ и $v_2$.
\begin{center}
\begin{tabularx}{\linewidth}{@{}p{0.31\linewidth}Y@{}}
\toprule
\textbf{Ситуация} & \textbf{Относительная скорость} \\
\midrule
Движение навстречу друг другу & $v_{\text{отн}}=v_1+v_2$.\\
Движение в одном направлении, быстрый догоняет медленного & $v_{\text{отн}}=v_{\text{быстр}}-v_{\text{медл}}$, причём $v_{\text{быстр}}>v_{\text{медл}}$.\\
Один объект неподвижен & $v_{\text{отн}}=v$ --- скорость движущегося объекта.\\
\bottomrule
\end{tabularx}
\end{center}
Если скорости при движении в одном направлении равны, расстояние между объектами остаётся неизменным: при ненулевом промежутке догоняния не произойдёт.

\section*{3. Единицы измерения: обязательная проверка}
Если скорость дана в км/ч, длины удобно выражать в километрах, а время --- в часах.
\[
\boxed{\displaystyle t_{\text{ч}}=\frac{t_{\text{с}}}{3600}},\qquad
\boxed{\displaystyle t_{\text{ч}}=\frac{t_{\text{мин}}}{60}}.
\]
\[
\boxed{\displaystyle S_{\text{км}}=\frac{S_{\text{м}}}{1000}},\qquad
\boxed{\displaystyle S_{\text{м}}=1000S_{\text{км}}}.
\]
Например, $18$ с $=18/3600$ ч, $1$ мин $=1/60$ ч, $900$ м $=0{,}9$ км. Перед составлением уравнения приведите \emph{все} данные к согласованным единицам.

\clearpage
\section*{4. Как определить относительный путь}
\begin{center}
\begin{tabularx}{\linewidth}{@{}p{0.45\linewidth}Y@{}}
\toprule
\textbf{Что происходит} & \textbf{Какой путь подставляем} \\
\midrule
Поезд полностью проходит мимо столба или человека & $S_{\text{отн}}=L_{\text{поезда}}$; длину человека принимаем равной нулю в масштабе задачи.\\
Поезд полностью проходит мост или лесополосу длиной $B$ & $S_{\text{отн}}=L_{\text{поезда}}+B$.\\
Два поезда полностью разъезжаются или один полностью обгоняет другой & $S_{\text{отн}}=L_1+L_2$.\\
\bottomrule
\end{tabularx}
\end{center}
\key{Уточнение:} формула с суммой длин относится к \emph{полному} прохождению объектов. Если в условии зафиксирован другой момент (например, только встреча передних частей), расстояние нужно определить специально.

\ex{Поезд проходит мимо столба}
Скорость поезда $60$ км/ч; прохождение длится $30$ с. Найдите длину поезда.
\[
L=60\cdot\frac{30}{3600}=0{,}5\ \text{км}=\boxed{500\ \text{м}}.
\]

\ex{Поезд и пешеход движутся навстречу}
Скорости $76$ км/ч и $4$ км/ч, время полного прохождения $18$ с. Пешеход считается точечным ориентиром:
\[
L=(76+4)\cdot\frac{18}{3600}=0{,}4\ \text{км}=\boxed{400\ \text{м}}.
\]
Здесь скорости \emph{складываются}, поскольку расстояние между встречными объектами сокращается быстрее.

\ex{Поезд проходит лесополосу}
Поезд со скоростью $72$ км/ч проходит лесополосу длиной $900$ м за $1$ мин. При $L$ --- длине поезда в километрах:
\[
\frac1{60}=\frac{L+0{,}9}{72}
\quad\Longrightarrow\quad
L=\frac{72}{60}-0{,}9=0{,}3\ \text{км}.
\]
\ans{Ответ: $300$ м.} От общей длины относительного пути вычли длину лесополосы.

\clearpage
\section*{5. Два движущихся поезда и точка отсчёта}
Если два состава полностью проходят мимо друг друга, учитываются \emph{обе длины}. При встречном движении их относительная скорость равна сумме скоростей.

\ex{Два поезда навстречу друг другу}
Длины составов $320$ и $280$ м, скорости $72$ и $36$ км/ч. Найдите время полного разъезда.
\[
S_{\text{отн}}=0{,}32+0{,}28=0{,}6\ \text{км},\qquad
v_{\text{отн}}=72+36=108\ \text{км/ч}.
\]
\[
t=\frac{0{,}6}{108}\ \text{ч}=\frac{0{,}6\cdot3600}{108}\ \text{с}
=\boxed{20\ \text{с}}.
\]

\key{Важная ловушка.} Человек находится в пассажирском поезде и измеряет время, за которое \emph{товарный поезд проходит мимо него}. В этой ситуации ориентир --- человек. Длина пассажирского поезда в расчёт относительного пути не входит, хотя его скорость влияет на относительную скорость.

\ex{Наблюдатель внутри движущегося поезда}
Пассажирский поезд движется со скоростью $60$ км/ч, навстречу ему --- товарный со скоростью $30$ км/ч. Товарный состав проходит мимо одного пассажира за $20$ с. Тогда
\[
L_{\text{тов}}=(60+30)\cdot\frac{20}{3600}
=0{,}5\ \text{км}=\boxed{500\ \text{м}}.
\]
Если бы спрашивалось время \emph{полного прохождения двух составов} мимо друг друга, понадобились бы длины обоих поездов.

\section*{6. Контроль ошибок при движении}
\begin{itemize}
\item \key{Неправильный знак:} встречное движение требует суммы скоростей; догоняние --- положительной разности.
\item \key{Неверный ориентир:} пассажир, столб и целый поезд задают разные относительные пути.
\item \key{Потеря длины моста:} при полном прохождении путь равен сумме длины поезда и препятствия.
\item \key{Несогласованные единицы:} секунды и метры нельзя непосредственно подставлять вместе с км/ч.
\item \key{Неправдоподобный результат:} длина, время и относительная скорость должны быть положительными.
\end{itemize}

\clearpage
\section*{7. Обгон с промежутками и отрывом}
Иногда сначала нужно \emph{догнать} другой объект, затем полностью его обогнать, после чего оторваться на заданное расстояние. При движении в одном направлении быстрый объект должен преодолеть относительно медленного:
\[
\boxed{\displaystyle S_{\text{отн}}=d_{\text{нач}}+L_1+L_2+d_{\text{кон}}.}
\]
Здесь $d_{\text{нач}}$ --- исходный чистый промежуток между ближайшими концами объектов; $d_{\text{кон}}$ --- требуемый чистый промежуток между ними после полного обгона. Тогда
\[
\boxed{\displaystyle v_{\text{быстр}}-v_{\text{медл}}
=\frac{d_{\text{нач}}+L_1+L_2+d_{\text{кон}}}{t}.}
\]

\ex{Две баржи}
Длина быстрой баржи $40$ м, медленной --- $60$ м. В начале расстояние от носа задней баржи до кормы передней составляет $200$ м. После обгона от кормы быстрой до носа медленной баржи должно быть $300$ м. На весь манёвр ушло $18$ мин. Насколько скорость быстрой баржи больше?
\[
S_{\text{отн}}=40+60+200+300=600\ \text{м}=0{,}6\ \text{км}.
\]
\[
t=\frac{18}{60}=0{,}3\ \text{ч},\qquad
\Delta v=\frac{0{,}6}{0{,}3}=\boxed{2\ \text{км/ч}}.
\]
\key{Обозначение:} $\Delta v$ («дельта $v$») означает разность скоростей. В этой задаче значение $\Delta v$ положительно.

\key{Как проверить модель:} общая относительная дистанция складывается из четырёх отрезков: начального промежутка, двух длин и конечного отрыва. Если один из отрезков отсутствует в условии, его не нужно добавлять произвольно.

\section*{8. Порядок решения задач о протяжённых телах}
\begin{enumerate}
\item Уточните, кто движется и относительно какого объекта измеряется время.
\item Установите нужный относительный путь: одна длина, сумма длин или сумма с дополнительными промежутками.
\item Выберите относительную скорость: сумма либо положительная разность.
\item Переведите время и расстояния в согласованные единицы.
\item Запишите $S_{\text{отн}}=v_{\text{отн}}t$, найдите неизвестное и проверьте единицы ответа.
\end{enumerate}

\clearpage
\thispagestyle{plain}
\begin{center}
{\Large\color{darkaccent} Алгоритм: движение протяжённых тел}\par
\end{center}
\vspace{10mm}
\begin{center}
\begin{tikzpicture}[
 node distance=12mm,
 boxstep/.style={draw=darkaccent,rounded corners=2pt,align=center,text width=90mm,minimum height=12mm,inner sep=4pt,font=\small},
 choice/.style={draw=darkaccent,diamond,aspect=2.25,align=center,text width=39mm,inner sep=2pt,font=\small},
 branch/.style={draw=darkaccent,rounded corners=2pt,align=center,text width=57mm,minimum height=15mm,inner sep=4pt,font=\small},
 terminal/.style={draw=darkaccent,ellipse,align=center,text width=50mm,minimum height=12mm,inner sep=3pt,font=\small},
 arrow/.style={-{Stealth[length=2.5mm]},line width=0.75pt,darkaccent}
]
\node[terminal] (start) at (0,0) {Начало};
\node[boxstep] (ref) at (0,-2.05) {Установить, кто кого \emph{полностью} проходит и что означает измеренное время};
\node[choice] (dir) at (0,-4.55) {Движение навстречу?};
\node[branch] (plus) at (-3.55,-7.0) {$v_{\text{отн}}=v_1+v_2$};
\node[branch] (minus) at (3.55,-7.0) {$v_{\text{отн}}=v_{\text{быстр}}-v_{\text{медл}}$\\при догонянии;\\для столба $v_{\text{медл}}=0$};
\node[boxstep] (dist) at (0,-9.5) {Найти $S_{\text{отн}}$: одна длина, сумма длин или путь с промежутками};
\node[boxstep] (units) at (0,-11.75) {Перевести метры в километры, минуты или секунды --- в часы};
\node[boxstep] (equ) at (0,-14) {Составить уравнение\\ $S_{\text{отн}}=v_{\text{отн}}t$};
\node[boxstep] (check) at (0,-16.25) {Найти неизвестное, проверить положительность и единицы};
\node[terminal] (end) at (0,-18.3) {Ответ};
\draw[arrow] (start)--(ref);
\draw[arrow] (ref)--(dir);
\draw[arrow] (dir.west) -| node[pos=0.25,above,font=\small] {да} (plus.north);
\draw[arrow] (dir.east) -| node[pos=0.25,above,font=\small] {нет} (minus.north);
\coordinate (join) at (0,-8.42);
\draw[darkaccent,line width=0.75pt] (plus.south) |- (join);
\draw[darkaccent,line width=0.75pt] (minus.south) |- (join);
\draw[arrow] (join) -- (dist.north);
\draw[arrow] (dist)--(units);
\draw[arrow] (units)--(equ);
\draw[arrow] (equ)--(check);
\draw[arrow] (check)--(end);
\end{tikzpicture}
\end{center}
\clearpage

\section*{9. Производительность: аналог скорости}
Производительность $p$ показывает, какую часть работы или какое количество продукции выполняют за единицу времени. Используем три взаимосвязанные величины:
\[
\boxed{\displaystyle A=pt},\qquad
\boxed{\displaystyle p=\frac At},\qquad
\boxed{\displaystyle t=\frac Ap}.
\]
$A$ --- объём выполненной работы (результат), $p$ --- производительность, $t$ --- время. Для \emph{одного полного заказа} удобно принять $A=1$; тогда производительность работающего в одиночку за $T$ часов равна $1/T$ заказа в час.

\key{Основное правило совместной работы:} если исполнители одновременно выполняют один заказ, их производительности складываются:
\[
\boxed{\displaystyle p_{\text{общ}}=p_1+p_2.}
\]
Это правило предполагает постоянную производительность каждого исполнителя и отсутствие взаимных помех.

\section*{10. Табличный метод: первый мастер, второй и вместе}
\ex{Один заказ; 3 часа, 6 часов и совместная работа}
Первый мастер выполняет заказ за $3$ ч, второй --- за $6$ ч. Пусть совместное время равно $x>0$ часов.
\begin{center}
\begin{tabularx}{\linewidth}{@{}Yccc@{}}
\toprule
\textbf{Исполнитель} & \textbf{Работа} & \textbf{Время, ч} & \textbf{Производит.} \\
\midrule
Первый & $1$ & $3$ & $1/3$ \\
Второй & $1$ & $6$ & $1/6$ \\
Вместе & $1$ & $x$ & $1/x$ \\
\bottomrule
\end{tabularx}
\end{center}
При совместной работе сумма индивидуальных производительностей равна общей:
\[
\frac1x=\frac13+\frac16=\frac12
\quad\Longrightarrow\quad \boxed{x=2\ \text{ч}}.
\]
Проверка: за $2$ ч первый выполнит $2/3$ заказа, второй $2/6=1/3$; в сумме получаем один целый заказ.

\key{Памятка для самостоятельного решения:} заполните три строки таблицы
«первый, второй, вместе»; получите $p_1+p_2=p_{\text{общ}}$ и проверьте,
что совместное время меньше каждого индивидуального.

\clearpage
\section*{Тренировочный блок}
\textit{Путь А --- 10 заданий. Упражнения расположены от базовых к более сложным. Ответы находятся на отдельной странице.}

\begin{spacing}{1.12}\small
\subsection*{Базовый уровень}
\task{1}{Поезд идёт со скоростью $72$ км/ч и полностью проходит мимо неподвижного дорожного знака за $25$ с. Найдите длину поезда в метрах.}

\task{2}{Один мастер выполняет заказ за $4$ ч, другой --- за $12$ ч. За сколько часов они выполнят этот заказ вместе при постоянной производительности?}

\task{3}{Поезд длиной $L$ м со скоростью $54$ км/ч полностью проходит мост длиной $450$ м за $50$ с. Найдите $L$.}

\subsection*{Средний уровень}
\task{4}{Поезд со скоростью $72$ км/ч полностью проходит мимо пешехода, движущегося \emph{навстречу} со скоростью $6$ км/ч, за $24$ с. Найдите длину поезда.}

\task{5}{Первый насос наполняет пустой резервуар за $5$ ч, второй --- за $10$ ч. За сколько \emph{минут} оба насоса наполнят его вместе?}

\task{6}{Поезд со скоростью $60$ км/ч полностью проходит лесополосу длиной $500$ м за $1$ мин. Найдите длину поезда.}

\task{7}{Поезд со скоростью $90$ км/ч обгоняет велосипедиста, едущего в том же направлении со скоростью $6$ км/ч. Поезд полностью проходит мимо велосипедиста за $18$ с. Найдите длину поезда.}

\subsection*{Повышенный уровень}
\task{8}{Один мастер выполняет заказ в одиночку за $6$ ч, а вместе со вторым --- за $4$ ч. За сколько часов второй мастер выполнит заказ в одиночку?}

\task{9}{Пассажирский поезд со скоростью $54$ км/ч идёт навстречу товарному со скоростью $36$ км/ч. Пассажир заметил, что товарный состав полностью прошёл \emph{мимо него} за $14$ с. Какова длина товарного состава?}

\task{10}{Две баржи движутся в одном направлении. Быстрая баржа имеет длину $40$ м, медленная --- $60$ м. Вначале от носа задней баржи до кормы передней было $150$ м. Через $15$ мин от начала манёвра, после полного обгона, расстояние от кормы быстрой баржи до носа медленной составило $250$ м. На сколько км/ч скорость быстрой баржи больше скорости медленной?}

\end{spacing}
\clearpage
\section*{Ответы к тренировочному блоку}
\begin{center}
\renewcommand{\arraystretch}{1.5}
\begin{tabularx}{0.89\linewidth}{@{}cY@{}}
\toprule
\textbf{№} & \textbf{Ответ} \\
\midrule
1 & $500$ м \\
2 & $3$ ч \\
3 & $300$ м \\
4 & $520$ м \\
5 & $200$ мин \\
6 & $500$ м \\
7 & $420$ м \\
8 & $12$ ч \\
9 & $350$ м \\
10 & $2$ км/ч \\
\bottomrule
\end{tabularx}
\end{center}

\vspace{8mm}
\section*{Итоговая памятка}
\begin{tabularx}{\linewidth}{@{}p{0.39\linewidth}Y@{}}
\toprule
\textbf{Сигнал в условии} & \textbf{Что делать} \\
\midrule
«Навстречу» & Сложить скорости. \\
«Догоняет», «обгоняет» & Вычесть меньшую скорость из большей. \\
«Мимо человека или столба» & Взять длину движущегося протяжённого тела. \\
«Полностью прошёл мост» & Сложить длину поезда и длину моста. \\
«Полностью прошли два поезда» & Сложить обе длины. \\
«Начальный промежуток, конечный отрыв» & Добавить оба промежутка к сумме длин. \\
«Вместе выполняют один заказ» & Сложить индивидуальные производительности. \\
«Скорость в км/ч» & Перевести метры в километры, секунды/минуты --- в часы. \\
\bottomrule
\end{tabularx}
\vfill
\begin{center}
\small\color{hint} Лёвин Артём Александрович --- эксклюзивно для Анны Банновой
\end{center}
\end{document}
